%% parts/p07/proof-complexity.tex --- Phase-4 P7 new-content slice.
%% Coverage: D-P8-007, N-c-027, A-075, and the P7 contributions to A-186/A-189.
%% The chapter states cited lower bounds; it contains no body proof.

\chapter{Propositional Proof Complexity}
\label{sec:p07-proof-complexity}
\chaptermark{Propositional Proof Complexity}

\section{Cook--Reckhow proof systems}
\label{sec:p07-cook-reckhow}

\begin{hbdefinition}{def:p07-cook-reckhow}{Cook--Reckhow propositional proof system}
Fix a self-delimiting binary encoding of propositional formulas, and let
\(\mathsf{TAUT}\) be the set of codes of classical tautologies.  A
\emph{Cook--Reckhow propositional proof system} is a polynomial-time
computable surjection
\[
  \mathsf P_{\mathrm{pf}}:\{0,1\}^{*}\longrightarrow\mathsf{TAUT}.
\]
Equivalently, after fixing a default tautology for malformed inputs, it is a
polynomial-time decidable relation
\(\operatorname{Prf}_{\mathsf P_{\mathrm{pf}}}(\pi,\varphi)\) whose range is exactly
\(\mathsf{TAUT}\).  For any proof or formula object \(a\) used in this chapter,
define \(\operatorname{size}(a)\) to be the bit length of its fixed encoding,
not a cardinality.

The name is necessary because efficient verification, rather than proof
search or one particular calculus, is the comparison invariant.  Its standing
is an abstract encoded proof system.  It differs from a calculus sheet by
forgetting presentation details while retaining a polynomial-time verifier,
and this nonredundancy is used in the boundedness theorem and simulation table
below \cite{CookReckhow1979}.
\index{Cook--Reckhow proof system}\index{proof size}
\end{hbdefinition}

\begin{hbdefinition}{def:p07-proof-simulation}{Boundedness and p-simulation}
A Cook--Reckhow system \(\mathsf P_{\mathrm{pf}}\) is \emph{polynomially bounded} iff
there is a polynomial \(q\) such that every tautology \(\varphi\) has a
\(\mathsf P_{\mathrm{pf}}\)-proof \(\pi\) with
\(\operatorname{size}(\pi)\leq q(\operatorname{size}(\varphi))\).
The system \(\mathsf P_{\mathrm{pf}}\) \emph{p-simulates}
\(\mathsf Q_{\mathrm{pf}}\) iff one polynomial-time transformation sends every
\(\mathsf Q_{\mathrm{pf}}\)-proof of \(\varphi\) to a
\(\mathsf P_{\mathrm{pf}}\)-proof of the same \(\varphi\), with
polynomially bounded output size.

These are metalinguistic properties of encoded proof systems.  A nonuniform
comparison of shortest proof lengths supplies no effective translation and
therefore is not p-simulation.  Table~\ref{tab:p07-proof-complexity-systems}
uses only the uniform notion.  Polynomial boundedness is instead an internal
short-proof-existence property of one system; p-simulation is a uniform
translation between two systems, and the theorem below uses the former.
\index{p-simulation}\index{polynomially bounded proof system}
\end{hbdefinition}

\begin{hbtheorem}{thm:p07-cook-reckhow-np-conp}{cited=CookReckhow1979}{Polynomially bounded proof systems and NP versus coNP}
A polynomially bounded Cook--Reckhow propositional proof system exists if and
only if \(\mathsf{NP}=\mathsf{coNP}\).  The theorem concerns existence of
some encoded system; it does not say that resolution, Frege, or extended
Frege is polynomially bounded \cite{CookReckhow1979,Krajicek2019}.
\end{hbtheorem}

\section{Resolution and the pigeonhole principle}
\label{sec:p07-resolution-lower-bound}

\begin{hbdefinition}{def:p07-pigeonhole-cnf}{The pigeonhole-principle CNFs}
For \(n\geq1\), let \(x_{i,j}\) say that pigeon \(i\) occupies hole \(j\),
where \(1\leq i\leq n+1\) and \(1\leq j\leq n\).  The clause set
\(\mathsf{PHP}^{n}_{n+1}\) consists of
\[
  x_{i,1}\lor\cdots\lor x_{i,n}
  \quad(1\leq i\leq n+1)
\]
and
\[
  \neg x_{i,j}\lor\neg x_{i',j}
  \quad(1\leq i<i'\leq n+1,\ 1\leq j\leq n).
\]
No clauses forbidding one pigeon from occupying two holes are included; that
stronger encoding must be named separately.  A resolution refutation uses
\hyperref[sheet:resolution-cpl]{\texttt{sheet:resolution-cpl}}, and its size
here is the number of clause occurrences in its derivation DAG, including
inputs.  Fixing the encoding and measure is necessary because lower bounds
for a different PHP presentation or proof-size measure are not literally the
same theorem; the present convention is used immediately below.
\index{pigeonhole principle!CNF}\index{resolution!proof size}
\end{hbdefinition}

\begin{hbtheorem}{thm:p07-haken-resolution-lower-bound}{sketched=annex:p7.complexity.1}{Exponential resolution lower bound for pigeonhole CNFs}
There is a real constant \(c>1\) such that, for all sufficiently large
\(n\), every resolution refutation of
\(\mathsf{PHP}^{n}_{n+1}\) has at least \(c^n\) clause occurrences.  This is
a lower bound for the fixed resolution system and encoding, not for Frege,
extended Frege, or arbitrary unsatisfiable CNFs
\cite{Haken1985,BeamePitassi1998}.
\end{hbtheorem}

\section{Frege and extension systems}
\label{sec:p07-frege-systems}

\begin{hbdefinition}{def:p07-frege-systems}{Frege and extended Frege systems}
In this chapter, a \emph{Frege system} is a fixed finite collection of sound
propositional axiom schemata and inference rules that is complete for
classical propositional logic and whose finite derivations are recognizable
in polynomial time.  The concrete handbook representative is
\hyperref[sheet:hilbert-cpl]{\texttt{sheet:hilbert-cpl}}.

An \emph{extended Frege system} is that representative together with the
fresh, acyclic definitional declarations on the
\hyperref[sheet:extended-frege-cpl]{extended-Frege sheet}.
An \emph{extended-resolution system} is propositional resolution together
with the fresh, acyclic three-clause definitional packages on
the \hyperref[sheet:extended-resolution-cpl]{extended-resolution sheet}.
These names are needed to distinguish formula-line proofs from clause
refutations and to distinguish mere abbreviation from unrestricted new
axioms.  Their standing is system-level, not semantic.  Freshness and
acyclicity make the extensions conservative abbreviations, so the extension
systems are not redundant names for their unextended bases; the distinction
is used in the simulation and open-problem cells of
Table~\ref{tab:p07-proof-complexity-systems}.
\index{Frege system}\index{extended Frege}\index{extended resolution}
\end{hbdefinition}

\begin{hbtheorem}{thm:p07-frege-lower-bound-boundary}{cited=Alekseev2026}{Frege lower bounds}
No superpolynomial lower bound is known for arbitrary-depth Frege proofs, and
none is known for extended Frege proofs.  Haken's resolution lower bound does
not transfer to either system.  This 2026 status statement is cited rather
than treated as a theorem of the handbook's formal development
\cite{Krajicek2019,Alekseev2026}.
\end{hbtheorem}

The source roles are specific.  Kraj\'i\v{c}ek supplies the proof-system
background and historical lower-bound boundary; Alekseev's 2026 introduction
explicitly identifies a superpolynomial arbitrary-depth Frege lower bound as a
major open problem and an extended-Frege lower bound as still harder.  The
paper's algebraic lower bound is not asserted to be either of those two
missing Boolean lower bounds.

\begin{landscape}
\thispagestyle{plain}
\section{Proof-system comparison}
\label{sec:p07-proof-system-comparison}
\begingroup\scriptsize
\setlength{\tabcolsep}{3pt}
\begin{longtable}{@{}p{.70in}p{1.70in}p{.80in}p{1.65in}p{1.80in}@{}}
\caption{Proof-complexity systems.  Every simulation claim
  uses Definition~\ref{def:p07-proof-simulation}; no shortest-proof-only
  comparison is substituted.}
\label{tab:p07-proof-complexity-systems}\\
\toprule
System & Proof objects and rule surface & Verification & Representative lower-bound status & Simulation statement \\
\midrule
\endfirsthead
\toprule
System & Proof objects and rule surface & Verification & Representative lower-bound status & Simulation statement \\
\midrule
\endhead
Resolution & Clause DAGs using \texttt{sheet:resolution-cpl} & Polynomial time & Exponential for the fixed \(\mathsf{PHP}^{n}_{n+1}\) family, Theorem~\ref{thm:p07-haken-resolution-lower-bound} & Baseline; no converse simulation by resolution is asserted. \\
\addlinespace
Extended resolution & Resolution plus the acyclic extension packages on \texttt{sheet:extended-resolution-cpl} & Polynomial time & No general superpolynomial lower bound asserted here & P-simulates ordinary resolution by using no extension declarations. \\
\addlinespace
Frege & Formula-line derivations in the fixed finite basis \texttt{sheet:hilbert-cpl} & Polynomial time & No general superpolynomial lower bound known & The fixed representative p-simulates the displayed resolution calculus after the encodings above are fixed.  Any two fixed finite sound and complete Frege bases with polynomial-time recognizable instances p-simulate one another under the same encoding convention. \\
\addlinespace
Extended Frege & Frege derivations plus fresh acyclic extension declarations & Polynomial time & No general superpolynomial lower bound known & P-simulates its underlying Frege system; no converse is asserted in this table. \\
\bottomrule
\end{longtable}
\endgroup
\end{landscape}

\section{Further reading and navigation}
\label{sec:p07-further-reading}

The following is the Part~7 annotated reading route:
\begin{itemize}
  \item Sequent calculi, normalization, and reductions.
        \par\parencite{Gentzen1934a,Gentzen1935b}
        \par\parencite{TroelstraSchwichtenberg2000}
  \item Lambda calculus and types.
        \par\parencite{Barendregt1984,GirardLafontTaylor1989}
  \item Ordinal analysis. \par\parencite{Takeuti1987,Pohlers2009}
  \item Functional interpretations.
        \par\parencite{AvigadFeferman1998,Kohlenbach2008}
  \item Proof complexity. \par\parencite{CookReckhow1979,Krajicek2019}
\end{itemize}
Each cited result in the preceding chapters retains its own status rather than
inheriting authority from this list.

\begingroup
\pagestyle{plain}
\begin{landscape}
\thispagestyle{plain}
\section{Body, sheet, and annex routes}
\label{sec:p07-navigation}

Table~\ref{tab:p07-body-annex-routes} routes every separately authored
status-bearing Part-7 result and the inherited result families on which those
additions depend.  Rule notation is read from the named registry sheet.  A
cited result has no annex item.

\begingroup\scriptsize
\setlength{\tabcolsep}{3pt}
\begin{longtable}{@{}p{.82in}p{2.05in}p{4.95in}@{}}
\caption{Part-7 body-to-sheet and body-to-annex routes.}
\label{tab:p07-body-annex-routes}\\
\toprule
Body topic & Exact rule-sheet route & Status destinations for every theorem-class entry \\
\midrule
\endfirsthead
\toprule
Body topic & Exact rule-sheet route & Status destinations for every theorem-class entry \\
\midrule
\endhead
Calculi and derivability &
\texttt{sheet:lk}, \texttt{sheet:lk-defneg}, \texttt{sheet:lk-eq},
\texttt{sheet:lj}, \texttt{sheet:nd-cpl}, \texttt{sheet:nd-fol}, and
\texttt{sheet:nj} &
\hyperref[thm:p07-lkd-soundness]{Defined-negation soundness}
--- \hyperref[annex:p7.calculi.1]{sketch CALCULI.1};
\hyperref[thm:p07-lk-presentation-equivalence]{presentation equivalence}
--- \hyperref[annex:p7.calculi.2]{sketch CALCULI.2};
\hyperref[thm:p07-lk-nd-equivalence]{LK--ND equivalence}
--- \hyperref[annex:p7.calculi.3]{proof CALCULI.3};
\hyperref[thm:p07-canonical-lk-completeness]{canonical-LK completeness}
--- \hyperref[annex:p7.calculi.4]{proof CALCULI.4}. \\
\addlinespace
Cut elimination and interpolation &
\texttt{sheet:lk-defneg} and \texttt{sheet:lj}; canonical LK and
LK-with-equality are expressly outside the reduction theorem &
\hyperref[lem:cut-mix-equipollence]{Cut/Mix equipollence}
--- source complete proof;
\hyperref[ch:cut-elimination]{cut elimination}
--- source proof sketch;
\hyperref[lem:cutfree-subformula]{cut-free subformula}
--- source complete proof;
\hyperref[sec:cut-elimination-02]{analytic existence and empty-sequent underivability}
--- source proof sketches;
\hyperref[thm:p07-midsequent]{midsequent} and
\hyperref[thm:p07-herbrand-midsequent]{Herbrand route}
--- \texttt{cited=Takeuti1987};
\hyperref[thm:p07-maehara-partition]{Maehara partition}
--- \hyperref[annex:p7.interp.1]{proof INTERP.1};
\hyperref[thm:p07-relational-cut-free-completeness]{relational cut-free completeness}
--- \hyperref[annex:p7.interp.2]{proof INTERP.2};
\hyperref[thm:craig-interpolation]{Craig}
--- \hyperref[annex:p7.interp.3]{proof INTERP.3}, including functions and equality;
\hyperref[thm:beth]{Beth}
--- \hyperref[annex:p7.interp.4]{proof INTERP.4}, including uniform finite support. \\
\addlinespace
Normalization and Curry--Howard &
\texttt{sheet:nj} and \texttt{sheet:stlc} &
\hyperref[sec:normalization-curry-howard-01]{NJ normalization}
--- source proof sketch;
\hyperref[sec:normalization-curry-howard-02]{STLC subject reduction and strong normalization}
--- source proof sketches. \\
\addlinespace
Lambda calculi and type theories &
\texttt{sheet:lambda-untyped}, \texttt{sheet:sk-combinatory}, and
\texttt{sheet:system-f}; the broader MLTT rule surface is
\texttt{sheet:mltt-intensional}, while the cited decision theorem covers only
its stated dependent-function/one-universe subcalculus &
\hyperref[thm:p07-church-rosser]{Church--Rosser}
--- \hyperref[annex:p7.lambda.1]{proof LAMBDA.1};
\hyperref[thm:p07-beta-normal-form-unique]{normal-form uniqueness}
--- \hyperref[annex:p7.lambda.2]{proof LAMBDA.2};
\hyperref[thm:p07-lambda-partial-recursive]{lambda/recursive equivalence}
--- \texttt{cited=Kleene1936};
\hyperref[thm:p07-bracket-abstraction]{bracket abstraction}
--- \texttt{cited=CurryFeys1958};
\hyperref[thm:p07-system-f-confluence]{System-F confluence}
--- \texttt{cited=GirardLafontTaylor1989};
\hyperref[thm:p07-system-f-strong-normalization]{System-F normalization}
--- \texttt{cited=Girard1972};
\hyperref[thm:p07-system-f-church-checking]{Church-style checking}
--- \texttt{cited=GirardLafontTaylor1989};
\hyperref[thm:p07-system-f-curry-undecidable]{Curry-style undecidability}
--- \texttt{cited=Wells1999};
\hyperref[thm:p07-mltt-normalization-decision]{matched MLTT decision result}
--- \texttt{cited=AbelCoquandDybjer2007}, expressly not a theorem about the
whole active sheet;
\hyperref[thm:p07-coc-strong-normalization]{CoC normalization}
--- \texttt{cited=CoquandHuet1988}. \\
\addlinespace
Arithmetic proof theory &
\texttt{sheet:omega-arithmetic} for the infinitary presentation; the finite
set-support presentation is derived from the propositional restriction of
\texttt{sheet:lk-defneg}; the full arithmetic analysis remains source-relative &
\hyperref[thm:p07-prop-cut-termination]{finite-set propositional Cut strategy}
--- \hyperref[annex:p7.arithmetic.1]{proof ARITHMETIC.1};
\hyperref[thm:p07-omega-cut-elimination]{omega-calculus cut elimination}
--- source-cited theorem;
\hyperref[sec:normalization-curry-howard-04]{Gentzen analysis}
--- \texttt{cited=Gentzen1936};
\hyperref[thm:p07-feferman-schuette]{Feferman--Schuette calibration}
--- \texttt{cited=Feferman1964}. \\
\addlinespace
Functional interpretations &
\texttt{sheet:system-t} and \texttt{sheet:lj}; the cited HA and
realizability metatheorems do not acquire local rule sheets &
\hyperref[thm:p07-dialectica-soundness]{Dialectica soundness} and
\hyperref[thm:p07-dialectica-consistency]{its consistency corollary}
--- \texttt{cited=AvigadFeferman1998};
\hyperref[thm:p07-number-realizability-soundness]{number-realizability soundness}
--- \texttt{cited=Kleene1945};
\hyperref[thm:p07-ect-characterization]{ECT characterization}
--- \texttt{cited=Troelstra1998};
\hyperref[thm:p07-markov-realizability-boundary]{Markov boundary},
\hyperref[thm:p07-ipc-disjunction-property]{IPC disjunction},
\hyperref[thm:p07-ha-disjunction-property]{HA disjunction}, and
\hyperref[thm:p07-ha-numerical-existence]{HA numerical existence}
--- \texttt{cited=TroelstraVanDalen1988};
\hyperref[thm:p07-lj-disjunction-last-rule]{cut-free LJ extraction}
--- \hyperref[annex:p7.functional.1]{proof FUNCTIONAL.1}. \\
\addlinespace
Epsilon calculus and generalized frameworks &
\texttt{sheet:epsilon}, \texttt{sheet:gs5-avron},
\texttt{sheet:labelled-k-negri}, \texttt{sheet:display-k-wansing}, and
\texttt{sheet:deep-s5-stouppa} &
\hyperref[thm:p07-first-epsilon]{first epsilon},
\hyperref[thm:p07-extended-first-epsilon]{extended first epsilon}, and
\hyperref[thm:p07-second-epsilon]{second epsilon}
--- \texttt{cited=MoserZach2006};
\hyperref[thm:p07-gs5-adequacy]{GS5 adequacy}
--- \texttt{cited=Avron1996};
\hyperref[thm:p07-labelled-k-adequacy]{labelled-K adequacy}
--- \texttt{cited=Negri2005};
\hyperref[thm:p07-display-k-adequacy]{display-K adequacy}
--- \texttt{cited=Wansing1994};
\hyperref[thm:p07-deep-s5-adequacy]{deep-S5 adequacy}
--- \texttt{cited=Stouppa2007}. \\
\addlinespace
Propositional proof complexity &
\texttt{sheet:resolution-cpl}, \texttt{sheet:extended-resolution-cpl},
\texttt{sheet:hilbert-cpl}, and \texttt{sheet:extended-frege-cpl} &
\hyperref[thm:p07-cook-reckhow-np-conp]{Cook--Reckhow equivalence}
--- \texttt{cited=CookReckhow1979};
\hyperref[thm:p07-haken-resolution-lower-bound]{Haken lower bound}
--- \hyperref[annex:p7.complexity.1]{sketch COMPLEXITY.1};
\hyperref[thm:p07-frege-lower-bound-boundary]{Frege lower-bound boundary}
--- \texttt{cited=Alekseev2026}. \\
\bottomrule
\end{longtable}
\endgroup
\end{landscape}
\endgroup

The resolution sheets connect forward to
Chapter~\ref{sec:p12-first-order-automation}, where unification, refinements,
equality handling, and practical search boundaries are treated explicitly.
